% Deutsche Parallelfassung zu why/motivation.tex

\chapter{Motivation}

\section{Warum ein Eisenbahnübergang kein Detail des Straßenkorridors ist}

Ein Straßenkorridor beeinflusst den verfügbaren Fahrweg, ein
Eisenbahnfahrzeug ist jedoch kinematisch durch sein Schienenpaar geführt. Das
Gleis schreibt daher den Krümmungsverlauf vor, den das Fahrzeug erfährt. Eine
Änderung von \(\kappa(s)\), seinen Ableitungen, der Überhöhung \(u(s)\) oder
der Geschwindigkeit überträgt sich auf Rad--Schiene-Interaktion,
Fahrzeugantwort, Fahrgastkomfort, Gleisbelastung, Verschleiß und
sicherheitsbezogene Bewertung. Übergangsentwurf ist folglich ein zentrales
Eisenbahnproblem und nicht bloß eine zweckmäßige Interpolation zwischen
Kartenelementen.

Man betrachte zwei Gesetze mit derselben Länge und denselben Randkrümmungen.
Eine Klothoide besitzt eine konstante Krümmungsableitung. Eine polynomiale oder
trigonometrische Halbwelle kann stattdessen ausgewählte Endbedingungen an
\(\kappa'(s)\) und \(\kappa''(s)\) erfüllen. Ähnlich erscheinende
Mittellinien können daher unterschiedliche dynamische Folgen tragen.
Umgekehrt kann ein glatteres Endverhalten mehr Länge oder eine andere
Verteilung der Freiheitsgrade verlangen. Die Klothoide ist wichtig, einfach
und häufig kompakt; diese Vorteile machen sie nicht zum universellen
physikalischen Optimum
\cite{BrustadDalmo2020Review,Koc2017SmoothedEnds}.

Dieser Vergleich begründete die ursprüngliche Forschungsrichtung der Arbeit.
Berlinish suchte eine gemeinsame funktionale Grammatik für solche
Alternativen, während axtranNew ihre zulässigen Kompositionen berechnen sollte.
AIM erweiterte die Frage später von \emph{Welcher Übergang lässt sich
berechnen?} zu \emph{Wie kann das resultierende Ingenieurwissen seine
Bedeutung bewahren und nutzbar werden?}

\section{Die Krümmungsänderung}

Die zu Beginn eingeführte Krümmungsänderung verschiebt die nachfolgenden
Positionen und ändert die Tangentenrichtung. Bleiben Überhöhung und
Geschwindigkeit unverändert, ändern sich auch fahrzeugbezogene Größen.
Nebenbedingungen, die vor der Änderung erfüllt waren, können danach verletzt
sein; der übliche Eisenbahnentwurf koppelt Krümmung, Überhöhung,
Geschwindigkeit und deren Änderungsraten
\cite{EN13803_2017,Klein1998Trassierung}.

Die sichtbare Kurve ist daher der Anfang der Folge, nicht ihr Ende. Die
Änderung muss gegenüber Gelände und benachbarten Objekten interpretiert, mit
dem bestehenden oder vermessenen Gleis verglichen, anhand anwendbarer Regeln
geprüft und als Alternative zur Begutachtung vorgelegt werden. Bei jedem
Schritt kann eine andere Repräsentation zweckmäßig sein. Nicht verloren gehen
darf die Beziehung zwischen Alignment, Änderung, resultierenden Zuständen,
Evidenz und Entscheidungsautorität.

\section{Drei leistungsfähige Sichten, drei notwendige Grenzen}

GIS, CAD und BIM nähern sich dieser Situation aus unterschiedlichen Stärken.
GIS ordnet das Alignment in einen umfassenden räumlichen Zusammenhang ein und
unterstützt interoperablen Zugriff, Verarbeitung, Visualisierung und
Erschließung raumbezogener Informationen \cite{OGCStandardsProgram2025}. Diese
Breite ist für Grundstücke, Umwelt, Netze und räumliche Analyse unverzichtbar;
der Ingenieur muss jedoch weiterhin bestimmen, welcher verfügbare Kontext für
die Alignment-Änderung anwendbar ist.

CAD konzentriert geometrische Konstruktion, Kommunikation und direkte
Bearbeitung \cite{KimKumarSnider2015Geometry}. Es macht die
Krümmungsänderung anschaulich und erzeugt die für Entwurf und Lieferung
benötigte Geometrie. Sein Geometriemodell begründet für sich genommen jedoch
weder die Provenienz einer Beobachtung noch die Autorität hinter einer Regel
oder die dauerhafte Identität, die mit einer anderen Repräsentation geteilt
wird.

BIM organisiert Informationen über gebaute Anlagen für den Austausch zwischen
Softwareanwendungen und Projektbeteiligten. IFC macht diese Rolle durch
Schemata und arbeitsablaufspezifische Modellsichten ausdrücklich
\cite{BuildingSMARTIFC4x3Scope}. Diese Strukturen können
Alignment-Informationen tragen, ohne dadurch dessen intrinsische Konstruktion
oder das gesamte zur Bewertung einer Änderung benötigte Wissen zu besitzen.

Es geht nicht darum, unter den drei Ansätzen einen Sieger zu bestimmen. Das
Alignment muss vielmehr zwischen ihren Sichten wechseln können, ohne lediglich
zu einem GIS-Feature, einer CAD-Kurve oder einem BIM-Container zu werden. Eine
Repräsentation ist gerade deshalb nützlich, weil sie für einen Zweck etwas
auswählt und ausdrückt. Diese Auswahl bedeutet zugleich, dass keine einzelne
Repräsentation stillschweigend zum gesamten Engineering Object werden darf.

\section{Fragmentiertes Wissen im Handlungsmoment}

Ein Teil des relevanten Wissens ist in Software kodiert. Anderes befindet sich
in Normen, Projektanforderungen, Messberichten, Berechnungsverfahren und
Entwurfsentscheidungen. Wieder anderes verbleibt bei Ingenieuren, die wissen,
welche Regel anwendbar und welche scheinbare Abweichung folgenreich ist.
Fragmentierung wird kostspielig, wenn eine Handlung diese Teile gemeinsam
benötigt. Die Forschung zur Bauinformatik führt diese beständige
Fragmentierung auf Spezialisierung, heterogene Werkzeuge und nur teilweise
interoperable Zusammenarbeit zurück \cite{Turk2020Interoperability}.

Für die Krümmungsänderung benötigt der Ingenieur mehr als Datenzugriff. Die
Software muss wissen, welches Objekt geändert wird, welcher Zustand
beabsichtigt ist, welche Beobachtungen einen anderen Zustand beschreiben,
welche Qualifikation die Regeln auswählt und welche Größen von der Änderung
abhängen. Zugleich muss sie die Grenze zwischen Berechnung und Autorisierung
bewahren.

Dies ist die praktische Motivation dafür, Ingenieurwissen anwendbar zu machen.
Anwendbares Wissen ist nicht einfach Text, der neben Geometrie gespeichert
wird. Es besitzt Geltungsbereich, Provenienz, Bedingungen und eine
verantwortete Bedeutung; es muss mit dem Zustand und der Operation verbunden
sein, für die es maßgeblich ist.

\section{Warum Krümmung das Problem sichtbar macht}

Ein Eisenbahn-Alignment wird häufig als Folge von Geraden, Kreisbögen und
Übergangsbögen eingeführt. Diese Beschreibung ist nützlich, doch Krümmung legt
die darunterliegende generative Beziehung offen. Entlang der intrinsischen
Bogenlänge \(s\) bestimmt die Krümmung die Änderung des Richtungswinkels,
\[
  \theta'(s)=\kappa(s),
\]
und durch Integration entstehen ebene Pose und Position. Änderungen von
Krümmung und Überhöhung gehen außerdem in ingenieurtechnische Bewertungen von
Fahrzeugreaktion und Komfort ein
\cite{Kufver2000RailwayAlignment,BrustadDalmo2020Review}.

Krümmung verbindet damit konstruktive Absicht mit geometrischer und dynamischer
Folge. Die resultierenden Koordinaten besitzen dennoch nicht die Identität des
Alignments, und eine gemessene Kurve offenbart nicht automatisch das
beabsichtigte Krümmungsgesetz. Die Unterscheidung zwischen intrinsischer
Konstruktion, metrischer Realisierung, physischer Realisierung und Beobachtung
ist keine philosophische Verzierung. Sie bestimmt, was nach der Änderung
sicher abgeleitet werden kann.

\section{Die Anforderung an AIM}

Eine brauchbare Antwort muss vier Anforderungen zugleich erfüllen:

\begin{itemize}
\item die Alignment-Identität über zweckgebundene Repräsentationen hinweg
  bewahren;
\item die konstruktiven und zustandsbezogenen Abhängigkeiten offenlegen, durch
  die Änderungen Folgen erhalten;
\item beabsichtigte, realisierte und beobachtete Informationen getrennt und
  nachvollziehbar halten;
\item Ingenieurwissen anwenden, ohne einem Dateiformat, einer
  Softwarekomponente oder einem numerischen Solver Entscheidungsautorität
  zuzuweisen.
\end{itemize}

Alignment-Based Information Modelling wird entwickelt, um diese Anforderungen
zu erfüllen. Sein Ausgangspunkt ist kein größeres Dateiformat, sondern eine
Alignment-spezifische Darstellung dessen, was das Engineering Object ist, wie
es messbaren und betrieblichen Zustand erhält, wie Evidenz sich darauf bezieht
und wie Wissen an der Bewertung mitwirkt.

\begin{summary}
Die durch eine Krümmungsänderung offengelegte Schwierigkeit besteht nicht
darin, dass Software die neue Kurve nicht zeichnen könnte. Sie besteht darin,
dass die ingenieurtechnische Bedeutung der Änderung Werkzeuge, Zustände,
Evidenz, Regeln und Autorität überquert. AIM setzt an, indem es diese
Unterscheidungen bewahrt und ihre Abhängigkeiten ausdrücklich macht.
\end{summary}
